% GENERATED FILE. Edit canonical lessons, then run npm run generate:books.
% canonical-source-hash: fnv1a64:f1f902868545aa7a
% canonical-lessons: FR-C468-sincere, FR-C468-menteur, FR-C468-naif, FR-C468-ruse, FR-C468-mentir

\chapter{Sincères et menteurs}
\label{ch:fr-sinceres-et-menteurs}
\hlchaptermodality{468}

\begin{chapteropening}
\textbf{By the end of this chapter:} \emph{I can say whether someone is honest, naive or cunning.}

Complete the last lesson of chapter 468: I can say whether someone is honest, naive or cunning.
\end{chapteropening}

\section[sincère]{sincère — sincere, honest}

\label{lesson:FR-C468-sincere}



\begin{quote}
Before the new one: say the French for clumsy, then the French for scatterbrained.
\end{quote}

\subsection*{You'll want to know: sincère}
\textbf{sincère} — \textquotedblleft{}sincere, honest\textquotedblright{}.

\begin{quote}
Elle est toujours sincère avec ses amis.
\end{quote}

\begin{quote}
\emph{She is always honest with her friends.}
\end{quote}

The noun is \emph{la sincérité}.

\begin{cousinweb}
\emph{sincère} is a learned borrowing from Latin \emph{sincērus} \textquotedblleft{}pure, clean, whole\textquotedblright{} (EtymDB; Lewis \& Short).
\end{cousinweb}

\begin{grammarlens}[title={What you've built}]
One more word for telling, explaining and describing at B1.
\end{grammarlens}

\subsection*{Your turn}
\begin{itemize}
\raggedright
  \item \emph{Say it:} \emph{sincère}
  \item \emph{Say it:} \emph{sincère}, once more
  \item \emph{Say it:} a sentence of your own with \emph{sincère}
  \item \emph{Recall:} say the French for clumsy, then the French for scatterbrained, then say \emph{sincère} again
\end{itemize}

\begin{tcolorbox}[breakable,colback=teal!4,colframe=teal!35!black,title={Before you move on}]
What does sincère mean? (\textquotedblleft{}Sincere, honest\textquotedblright{}.) Say it once more.
\end{tcolorbox}
\section[le menteur]{le menteur — liar; lying}

\label{lesson:FR-C468-menteur}



\begin{quote}
Before the new one: say the French for scatterbrained, then the French for sincere.
\end{quote}

\subsection*{You'll want to know: le menteur}
\textbf{le menteur} — \textquotedblleft{}liar; lying\textquotedblright{}.

\begin{quote}
Ne le crois pas, c'est un menteur.
\end{quote}

\begin{quote}
\emph{Do not believe him, he is a liar.}
\end{quote}

A woman is \emph{la menteuse}. A lie is \emph{un mensonge}.

\begin{cousinweb}
\emph{le menteur} is inherited from Old French \emph{menteor} (EtymDB), built on \emph{mentir} \textquotedblleft{}to lie\textquotedblright{}, from Latin \emph{mentīrī} \textquotedblleft{}to lie\textquotedblright{} (Lewis \& Short).
\end{cousinweb}

\begin{grammarlens}[title={What you've built}]
One more word for telling, explaining and describing at B1.
\end{grammarlens}

\subsection*{Your turn}
\begin{itemize}
\raggedright
  \item \emph{Say it:} \emph{le menteur}
  \item \emph{Say it:} \emph{le menteur}, once more
  \item \emph{Say it:} a sentence of your own with \emph{le menteur}
  \item \emph{Recall:} say the French for scatterbrained, then the French for sincere, then say \emph{le menteur} again
\end{itemize}

\begin{tcolorbox}[breakable,colback=teal!4,colframe=teal!35!black,title={Before you move on}]
What does le menteur mean? (\textquotedblleft{}Liar; lying\textquotedblright{}.) Say it once more.
\end{tcolorbox}
\section[naïf]{naïf — naive}

\label{lesson:FR-C468-naif}



\begin{quote}
Before the new one: say the French for sincere, then the French for liar.
\end{quote}

\subsection*{You'll want to know: naïf}
\textbf{naïf} — \textquotedblleft{}naive\textquotedblright{}.

\begin{quote}
J'étais si naïf que je croyais tout.
\end{quote}

\begin{quote}
\emph{I was so naive that I believed everything.}
\end{quote}

Feminine \emph{naïve}. The noun is \emph{la naïveté}.

\begin{cousinweb}
\emph{naïf} is inherited from Latin \emph{nātīvus} \textquotedblleft{}inborn, natural\textquotedblright{} (EtymDB), from \emph{nāscī} \textquotedblleft{}to be born\textquotedblright{} (Lewis \& Short); the sense moved from \textquotedblleft{}natural\textquotedblright{} to \textquotedblleft{}unspoiled\textquotedblright{} and then to \textquotedblleft{}naive\textquotedblright{}.
\end{cousinweb}

\begin{grammarlens}[title={What you've built}]
One more word for telling, explaining and describing at B1.
\end{grammarlens}

\subsection*{Your turn}
\begin{itemize}
\raggedright
  \item \emph{Say it:} \emph{naïf}
  \item \emph{Say it:} \emph{naïf}, once more
  \item \emph{Say it:} a sentence of your own with \emph{naïf}
  \item \emph{Recall:} say the French for sincere, then the French for liar, then say \emph{naïf} again
\end{itemize}

\begin{tcolorbox}[breakable,colback=teal!4,colframe=teal!35!black,title={Before you move on}]
What does naïf mean? (\textquotedblleft{}Naive\textquotedblright{}.) Say it once more.
\end{tcolorbox}
\section[rusé]{rusé — cunning, crafty}

\label{lesson:FR-C468-ruse}



\begin{quote}
Before the new one: say the French for liar, then the French for naive.
\end{quote}

\subsection*{You'll want to know: rusé}
\textbf{rusé} — \textquotedblleft{}cunning, crafty\textquotedblright{}.

\begin{quote}
Le renard est rusé : il trouve toujours une sortie.
\end{quote}

\begin{quote}
\emph{The fox is cunning: it always finds a way out.}
\end{quote}

The noun \emph{une ruse} is a trick.

\begin{cousinweb}
\emph{rusé} comes from \emph{ruser}, first said of a hunted animal that doubles back. \emph{Ruser} is inherited from Latin \emph{recūsāre} \textquotedblleft{}to refuse, to draw back\textquotedblright{} (Lewis \& Short), itself built on \emph{causa}.
\end{cousinweb}

\begin{grammarlens}[title={What you've built}]
One more word for telling, explaining and describing at B1.
\end{grammarlens}

\subsection*{Your turn}
\begin{itemize}
\raggedright
  \item \emph{Say it:} \emph{rusé}
  \item \emph{Say it:} \emph{rusé}, once more
  \item \emph{Say it:} a sentence of your own with \emph{rusé}
  \item \emph{Recall:} say the French for liar, then the French for naive, then say \emph{rusé} again
\end{itemize}

\begin{tcolorbox}[breakable,colback=teal!4,colframe=teal!35!black,title={Before you move on}]
What does rusé mean? (\textquotedblleft{}Cunning, crafty\textquotedblright{}.) Say it once more.
\end{tcolorbox}
\section[mentir]{mentir — to lie, to tell lies}

\label{lesson:FR-C468-mentir}



\begin{quote}
Before the new one: say the French for naive, then the French for cunning.
\end{quote}

\subsection*{You'll want to know: mentir}
\textbf{mentir} — \textquotedblleft{}to lie, to tell lies\textquotedblright{}.

\begin{quote}
Il a menti à ses parents sur ses notes.
\end{quote}

\begin{quote}
\emph{He lied to his parents about his marks.}
\end{quote}

Present: \emph{je mens, il ment, nous mentons}. Past: \emph{j'ai menti}.

\begin{cousinweb}
\emph{mentir} is inherited, through Old French, from Vulgar Latin \emph{mentīre}, a form of Latin \emph{mentīrī} \textquotedblleft{}to lie\textquotedblright{} (EtymDB; Lewis \& Short).
\end{cousinweb}

\begin{grammarlens}[title={What you've built}]
One more word for telling, explaining and describing at B1.
\end{grammarlens}

\subsection*{Your turn}
\begin{itemize}
\raggedright
  \item \emph{Say it:} \emph{mentir}
  \item \emph{Say it:} \emph{mentir}, once more
  \item \emph{Say it:} a sentence of your own with \emph{mentir}
  \item \emph{Recall:} say the French for naive, then the French for cunning, then say \emph{mentir} again
\end{itemize}

\begin{tcolorbox}[breakable,colback=teal!4,colframe=teal!35!black,title={Before you move on}]
What does mentir mean? (\textquotedblleft{}To lie, to tell lies\textquotedblright{}.) Say it once more.
\end{tcolorbox}
